\begin{explanationbox}
	\textbf{\textcolor{CExplanation}{一句话直觉}}

	把图像在x轴下面的部分翻到上面去，就像将湖面倒影翻到真实世界一样，上下对称。

	\medskip
	\textbf{\textcolor{CExplanation}{核心拆解}}
	\begin{description}[style=nextline, leftmargin=2.8em, labelsep=0.5em, itemsep=0.45em]
		\item[\textbf{要点一（轴上方保留）：}] 对于原函数中 $f(x)\ge0$ 的部分，保持不变，直接保留原图像。
		\item[\textbf{要点二（轴下方翻折）：}] 对于原函数中 $f(x)<0$ 的部分，将图像关于 $x$ 轴对称翻折到上方，即函数值取相反数。
	\end{description}

	\medskip
	\textbf{\textcolor{CExplanation}{几何本质（上下对称）}}

	翻折操作本质是将图形关于 $x$ 轴作对称，但只针对位于 $x$ 轴下方的部分。

	\medskip
	\textbf{\textcolor{CExplanation}{代数意义（分段函数）}}

	翻折后函数值等于原函数值的绝对值，因此可以写成分段形式，便于代数运算。

	\medskip
	\textbf{\textcolor{CExplanation}{考点价值（画图与求域）}}
	\begin{itemize}[leftmargin=*, itemsep=0.5em, topsep=0.4em]
		\item \textbf{考法一（快速画图）：} 给定 $y=f(x)$ 的图像，可直接通过翻折得到 $y=|f(x)|$ 的图像，用于后续分析。
		\item \textbf{考法二（值域与最值）：} 翻折后图像全部位于 $x$ 轴上方，值域是原函数值域的绝对值版本（负数变正数），便于求最值和比较函数值。
	\end{itemize}

	\medskip
	\textbf{\textcolor{CExplanation}{顿悟点}}

	翻折只改变函数值的符号，不改变自变量的取值范围，且翻折后函数值非负。

	\medskip
	\textbf{\textcolor{CExplanation}{使用场景}}
	\begin{itemize}[leftmargin=*, itemsep=0.5em, topsep=0.4em]
		\item \textbf{场景一（含绝对值函数作图）：} 当题目需要分析含绝对值的函数时，先画出原函数再翻折即可。
		\item \textbf{场景二（绝对值不等式直观化）：} 借助图像与 $y=k$ 交点求解 $|f(x)|\le k$ 等不等式。
	\end{itemize}
\end{explanationbox}
